\documentclass[10pt,a4paper]{article}
\usepackage{amsmath,amssymb,amsthm,graphicx}
% Card-local minimal preamble for S7 cards cut from sources/sigma-2025-001-symplectic-differential-reduction-algebras-a/pinned-original-source.tex
% (doi:10.3842/SIGMA.2025.001, SIGMA 21 (2025), CC BY 4.0). Author text untouched.
% Omitted and why: the publisher sigma document class (unavailable in the toolchain;
% the assembler supplies article), the enumitem/romanbar packages (not used inside
% the counted spans; romanbar is also absent from the host texlive). The theorem
% environments below are copied verbatim from the source preamble (source lines
% 11-19) with their shared per-section counter, the per-section equation numbering
% matches the source (source line 9), and the author macros used inside the counted
% spans (source lines 22-32) are copied verbatim; a qed stub is provided for the
% author proof-closing pattern.
\usepackage{amsmath,amssymb,amsthm,mathtools}
\usepackage{mathrsfs}
\usepackage{hyperref}
\newtheorem{Theorem}{Theorem}[section]
\newtheorem{Corollary}[Theorem]{Corollary}
\newtheorem{Lemma}[Theorem]{Lemma}
\newtheorem{Proposition}[Theorem]{Proposition}
 { \theoremstyle{definition}
\newtheorem{Definition}[Theorem]{Definition}
\newtheorem{Note}[Theorem]{Note}
\newtheorem{Example}[Theorem]{Example}
\newtheorem{Remark}[Theorem]{Remark} }
\numberwithin{equation}{section}
\DeclareMathOperator{\Id}{Id}
\DeclareMathOperator{\Span}{Span}
\DeclareMathOperator{\Ann}{Ann}
\DeclareMathOperator{\End}{End}
\DeclareMathOperator{\Hom}{Hom}
\DeclareMathOperator{\Aut}{Aut}
\newcommand{\II}{\mathrm{I\hspace{-.1em}I}}
\providecommand{\qed}{\ensuremath{\square}}
\usepackage{cleveref}
\crefname{originalproposition}{original proposition}{original propositions}
\Crefname{originalproposition}{Original proposition}{Original propositions}
\makeatletter
\newlabel{Prop:Normalized-generators-relations}{{4.1}{source}{}{sourcepointer-1}{}}
\newlabel{Prop:Normalized-generators-relations@cref}{{[originalproposition][0][]4.1}{[1][0][]source}{}{}{}}
\makeatother
\crefname{originalproposition}{original proposition}{original propositions}
\Crefname{originalproposition}{Original proposition}{Original propositions}
\makeatletter
\newlabel{Prop:Automorphisms}{{4.4}{source}{}{sourcepointer-2}{}}
\newlabel{Prop:Automorphisms@cref}{{[originalproposition][0][]4.4}{[1][0][]source}{}{}{}}
\makeatother
\crefname{originalequation}{original equation}{original equations}
\Crefname{originalequation}{Original equation}{Original equations}
\makeatletter
\newlabel{eq:GWA-rels}{{4.2}{source}{}{sourcepointer-3}{}}
\newlabel{eq:GWA-rels@cref}{{[originalequation][0][]4.2}{[1][0][]source}{}{}{}}
\makeatother
\crefname{originalequation}{original equation}{original equations}
\Crefname{originalequation}{Original equation}{Original equations}
\makeatletter
\newlabel{eq:GWA-rels1}{{4.2a}{source}{}{sourcepointer-4}{}}
\newlabel{eq:GWA-rels1@cref}{{[originalequation][0][]4.2a}{[1][0][]source}{}{}{}}
\makeatother
\crefname{originalequation}{original equation}{original equations}
\Crefname{originalequation}{Original equation}{Original equations}
\makeatletter
\newlabel{eq:GWA-rels2}{{4.2b}{source}{}{sourcepointer-5}{}}
\newlabel{eq:GWA-rels2@cref}{{[originalequation][0][]4.2b}{[1][0][]source}{}{}{}}
\makeatother
\crefname{originalequation}{original equation}{original equations}
\Crefname{originalequation}{Original equation}{Original equations}
\makeatletter
\newlabel{eq:rel-x1-H}{{3.7a}{source}{}{sourcepointer-6}{}}
\newlabel{eq:rel-x1-H@cref}{{[originalequation][0][]3.7a}{[1][0][]source}{}{}{}}
\makeatother
\crefname{originalequation}{original equation}{original equations}
\Crefname{originalequation}{Original equation}{Original equations}
\makeatletter
\newlabel{eq:rel-p2-H}{{3.7b}{source}{}{sourcepointer-7}{}}
\newlabel{eq:rel-p2-H@cref}{{[originalequation][0][]3.7b}{[1][0][]source}{}{}{}}
\makeatother
\crefname{originalequation}{original equation}{original equations}
\Crefname{originalequation}{Original equation}{Original equations}
\makeatletter
\newlabel{eq:rel-x1p1}{{3.7e}{source}{}{sourcepointer-8}{}}
\newlabel{eq:rel-x1p1@cref}{{[originalequation][0][]3.7e}{[1][0][]source}{}{}{}}
\makeatother
\crefname{originalcorollary}{original corollary}{original corollarys}
\Crefname{originalcorollary}{Original corollary}{Original corollarys}
\makeatletter
\newlabel{cor:Dsp4-structure:domain}{{3.6}{source}{}{sourcepointer-9}{}}
\newlabel{cor:Dsp4-structure:domain@cref}{{[originalcorollary][0][]3.6}{[1][0][]source}{}{}{}}
\makeatother
\crefname{originallemma}{original lemma}{original lemmas}
\Crefname{originallemma}{Original lemma}{Original lemmas}
\makeatletter
\newlabel{lemma:finite-generators}{{3.3}{source}{}{sourcepointer-10}{}}
\newlabel{lemma:finite-generators@cref}{{[originallemma][0][]3.3}{[1][0][]source}{}{}{}}
\makeatother
\title{D(sp(4)) is a generalized Weyl algebra: an explicit C-algebra isomorphism from B(sigma,t) onto D(sp(4)) in the normalized generators}
\author{Hartwig, Jonas T., Williams II, Dwight Anderson}
\date{doi:10.3842/SIGMA.2025.001; CC BY 4.0}
\begin{document}
\maketitle
\setcounter{section}{4}
\setcounter{Theorem}{5}
\noindent\textit{Editorial note.} All delivered mathematical text is the author's own, copied from the pinned source file sources/sigma-2025-001-symplectic-differential-reduction-algebras-a/pinned-original-source.tex; no statement or proof step was rewritten, reordered or removed. The counted span is the author environment and its own proof environment. References inside the span are handed to the assembler as navigation stubs; whether each of them is load-bearing is recorded on the stub itself.\par\medskip
\noindent\textit{Editorial note.} The in-place proof environment (source lines 595-597) is the author's own pointer to Appendix A and is not itself the counted proof; the counted proof span is the author's remote named proof environment at source lines 621-686 (021 tvbo remote named proof precedent form; no section command and no figure inside the block). The single citation inside the remote proof (bavula1992generalized, proof line 684) is resolved by a card-local bibliography fragment copied verbatim from the source thebibliography (source lines 712-714).\par\medskip
\noindent\textit{Editorial scope.} Single theorem card cut from doi:10.3842/SIGMA.2025.001 (Symplectic Differential Reduction Algebras and Generalized Weyl Algebras); statement and proof are the author own text at source lines 583-594 and 621-686.\par\medskip
\noindent\textit{Allowed standard prerequisites (editorial).}\begin{itemize}
\item Normalized generators (Proposition Prop:Normalized-generators-relations, printed 4.1). cited inside the counted statement and proof spans. Supplies the normalized generator relations that the GWA verification builds on.
\item Automorphisms (Proposition Prop:Automorphisms, printed 4.4). cited inside the counted statement and proof spans. Supplies the automorphism framework in which the normalized generators live.
\item GWA relations (equation eq:GWA-rels, printed 4.2). cited inside the counted proof span. The relation set verified by the extended argument; its two halves are eq:GWA-rels1 (4.2a) and eq:GWA-rels2 (4.2b).
\item Presentation relations (equations eq:rel-x1-H 3.7a through eq:rel-p2-H 3.7b, eq:rel-x1p1 3.7e). cited inside the counted proof span. The finite-presentation relations used to verify the H\_alpha, H\_beta cases of the GWA relations.
\item Domain corollary (Corollary cor:Dsp4-structure:domain, printed 3.6). cited inside the counted proof span. Yields that sigma\_1 and sigma\_2 commute in the final step of the extended argument.
\item Generators lemma (Lemma lemma:finite-generators, printed 3.3). cited inside the counted proof span. Supplies the left R-basis of D(sp(4)) that proves phi bijective.
\end{itemize}\medskip
\ifdefined\hypertarget\hypertarget{sourcepointer-1}{}\fi
\noindent\textit{Source reference.} Original proposition 4.1: sources/sigma-2025-001-symplectic-differential-reduction-algebras-a/pinned-original-source.tex, source lines 526--526. navigation stub only: the target lives outside the counted span in the same article; supplies the normalized generator relations the GWA verification builds on\par\medskip
\ifdefined\hypertarget\hypertarget{sourcepointer-2}{}\fi
\noindent\textit{Source reference.} Original proposition 4.4: sources/sigma-2025-001-symplectic-differential-reduction-algebras-a/pinned-original-source.tex, source lines 559--559. navigation stub only: the target lives outside the counted span in the same article; supplies the automorphism framework of the normalized generators\par\medskip
\ifdefined\hypertarget\hypertarget{sourcepointer-3}{}\fi
\noindent\textit{Source reference.} Original equation 4.2: sources/sigma-2025-001-symplectic-differential-reduction-algebras-a/pinned-original-source.tex, source lines 543--543. navigation stub only: the target lives outside the counted span in the same article; the GWA relation set verified by the extended argument\par\medskip
\ifdefined\hypertarget\hypertarget{sourcepointer-4}{}\fi
\noindent\textit{Source reference.} Original equation 4.2a: sources/sigma-2025-001-symplectic-differential-reduction-algebras-a/pinned-original-source.tex, source lines 545--545. navigation stub only: the target lives outside the counted span in the same article; first half of the GWA relations, reduced to generator cases in the proof\par\medskip
\ifdefined\hypertarget\hypertarget{sourcepointer-5}{}\fi
\noindent\textit{Source reference.} Original equation 4.2b: sources/sigma-2025-001-symplectic-differential-reduction-algebras-a/pinned-original-source.tex, source lines 547--547. navigation stub only: the target lives outside the counted span in the same article; second half of the GWA relations, verified by full computation in the proof\par\medskip
\ifdefined\hypertarget\hypertarget{sourcepointer-6}{}\fi
\noindent\textit{Source reference.} Original equation 3.7a: sources/sigma-2025-001-symplectic-differential-reduction-algebras-a/pinned-original-source.tex, source lines 401--401. navigation stub only: the target lives outside the counted span in the same article; presentation relation used for the H\_alpha and H\_beta cases\par\medskip
\ifdefined\hypertarget\hypertarget{sourcepointer-7}{}\fi
\noindent\textit{Source reference.} Original equation 3.7b: sources/sigma-2025-001-symplectic-differential-reduction-algebras-a/pinned-original-source.tex, source lines 410--410. navigation stub only: the target lives outside the counted span in the same article; presentation relation used for the H\_alpha and H\_beta cases\par\medskip
\ifdefined\hypertarget\hypertarget{sourcepointer-8}{}\fi
\noindent\textit{Source reference.} Original equation 3.7e: sources/sigma-2025-001-symplectic-differential-reduction-algebras-a/pinned-original-source.tex, source lines 425--425. navigation stub only: the target lives outside the counted span in the same article; presentation relation closing the pf-phi computation for the t\_i cases\par\medskip
\ifdefined\hypertarget\hypertarget{sourcepointer-9}{}\fi
\noindent\textit{Source reference.} Original corollary 3.6: sources/sigma-2025-001-symplectic-differential-reduction-algebras-a/pinned-original-source.tex, source lines 502--502. navigation stub only: the target lives outside the counted span in the same article; yields the sigma\_1, sigma\_2 commutation in the final step\par\medskip
\ifdefined\hypertarget\hypertarget{sourcepointer-10}{}\fi
\noindent\textit{Source reference.} Original lemma 3.3: sources/sigma-2025-001-symplectic-differential-reduction-algebras-a/pinned-original-source.tex, source lines 330--330. navigation stub only: the target lives outside the counted span in the same article; supplies the left R-basis of D(sp(4)) proving phi bijective\par\medskip
\begin{Theorem}[$D(\mathfrak{sp}(4))$ is a GWA]\label{thm:GWA-realization}
	With $B$ and $\sigma_{i}$ as in Proposition {\rm\ref{Prop:Normalized-generators-relations}}, there is a $\mathbb{C}$-algebra isomorphism
$\phi\colon B(\sigma,t) \longrightarrow D(\mathfrak{sp}(4))
$
	satisfying
	\[\phi(H_\alpha)=H_\alpha,\qquad
	\phi(H_\beta)=H_\beta,\qquad
	\phi(X_i)=\hat x_i,\qquad
	\phi(Y_i)=\hat \partial_i,\qquad i=1,2,
	\]
	where we use the normalized generators of $D(\mathfrak{sp}(4))$ from Proposition {\rm\ref{Prop:Automorphisms}}.
\end{Theorem}

\begin{proof}[Extended argument for proof of Theorem \ref{thm:GWA-realization}]
	Let $\phi\colon B\to D(\mathfrak{sp}(4))$ be the $\mathbb{C}$-algebra homomorphism determined by $\phi|_R=\Id_R$ and $\phi(t_i)=\hat\partial_i\hat x_i$ for $i=1,2$. Extend $\phi$ to $B\cup\{X_1,Y_1,X_2,Y_2\}$ by $\phi(X_i)=\hat x_i$, $\phi(Y_i)=\hat\partial_i$. We must show that the GWA relations \eqref{eq:GWA-rels} are satisfied by the images of $\phi$.
	First, we prove that relations \eqref{eq:GWA-rels2} are preserved.
	
	Using $\bar{x}_1 H_\alpha = (H_\alpha-1)\bar{x}_1$,
	\[
	\hat{x}_1\hat{x}_2 = \bar{x}_1(H_\alpha+2)\bar{x}_2 = (H_\alpha+1)\bar{x}_1\bar{x}_2=(H_\alpha+1)\frac{H_\alpha+2}{H_\alpha+1}\bar{x}_2\bar{x}_1=\hat{x}_2\hat{x}_1.
	\]
	
	Using $H_{\beta+\alpha}\bar{\partial}_1=\bar{\partial}_1(H_{\beta+\alpha}-1)$, $\bar{\partial}_2 H_\alpha=(H_\alpha-1)\bar{\partial}_2$, and $\bar{\partial}_2 H_{\beta+\alpha} = (H_{\beta+\alpha}+1)\bar{\partial}_2$, we get
	\begin{align*}
		\hat{\partial}_2\hat{\partial}_1 &= \bar{\partial}_2(H_{\beta+\alpha}+1)\bar{\partial}_1(H_\alpha+1)(H_{\beta+\alpha}+1)=\bar{\partial}_2\bar{\partial}_1(H_\alpha+1)H_{\beta+\alpha}(H_{\beta+\alpha}+1)\\
		&=\bar{\partial}_1\bar{\partial}_2\frac{H_\alpha+2}{H_\alpha+1}(H_\alpha+1)H_{\beta+\alpha}(H_{\beta+\alpha}+1)=
		\bar{\partial}_1\bar{\partial}_2(H_\alpha+2)H_{\beta+\alpha}(H_{\beta+\alpha}+1)\\
		&=\bar{\partial}_1 (H_\alpha+1)(H_{\beta+\alpha}+1)\bar{\partial}_2(H_{\beta+\alpha}+1)=\hat{\partial}_1\hat{\partial}_2.
	\end{align*}
	Next, since $H_{\beta+\alpha}$ commutes with $\bar{x}_1\bar{\partial}_2$, and $H_{\beta+\alpha}\bar{\partial}_2=\bar{\partial}_2(H_{\beta+\alpha}-1)$, we have
	\begin{align*}
		\hat{x}_1\hat{\partial}_2 &= \bar{x}_1\bar{\partial}_2(H_{\beta+\alpha}+1)=
		(H_{\beta+\alpha}+1)\bar{x}_1\bar{\partial}_2=
		(H_{\beta+\alpha}+1)\frac{H_{\beta+\alpha}+2}{H_{\beta+\alpha}+1}\bar{\partial}_2\bar{x}_1\\
		&=(H_{\beta+\alpha}+2)\bar{\partial}_2\bar{x}_1=\bar{\partial}_2(H_{\beta+\alpha}+1)\bar{x}_1=\hat{\partial}_2\hat{x}_1.
	\end{align*}
	From $H_\alpha\bar{x}_2=\bar{x}_2(H_\alpha-1)$, $H_\alpha\bar{\partial}_1=\bar{\partial}_1(H_\alpha-1)$, and $\bar{x}_2 H_{\beta+\alpha} = (H_{\beta+\alpha}-1)\bar{x}_2$, we see
	\begin{align*}
		\hat{x}_2\hat{\partial}_1&=(H_\alpha+2)\bar{x}_2\bar{\partial}_1(H_\alpha+1)(H_{\beta+\alpha}+1)=\bar{x}_2\bar{\partial}_1 H_\alpha (H_\alpha+1)(H_{\beta+\alpha}+1)\\
		&=\bar{\partial}_1\bar{x}_2\frac{H_{\beta+\alpha}+2}{H_{\beta+\alpha}+1}H_\alpha (H_\alpha+1)(H_{\beta+\alpha}+1)=
		\bar{\partial}_1\bar{x}_2 H_\alpha(H_\alpha+1)(H_{\beta+\alpha}+2)\\
		&=\bar{\partial}_1(H_\alpha+1)(H_\alpha+2)(H_{\beta+\alpha}+1)\bar{x}_2=\hat{\partial}_1\hat{x}_2.
	\end{align*}
	We have proved that the four relations in \eqref{eq:GWA-rels2} are preserved by $\phi$.
	
	Next, to show that \eqref{eq:GWA-rels1} holds for all $b\in B$, it suffices to show \eqref{eq:GWA-rels1} holds for $b\in\{H_\alpha, H_\beta, t_1, t_2\}$ since $\sigma_i$ are $\mathbb{C}$-algebra homomorphisms.
	For $b=H_\alpha$ and $b=H_\beta$, \eqref{eq:GWA-rels1} follows directly from \eqref{eq:rel-x1-H}--\eqref{eq:rel-p2-H} and the definition of $\sigma_i$.
	
	For $b=t_i$, the identities are immediately verified once we prove that $\phi(t_i)=\phi(Y_i)\phi(X_i)$ and~${\phi\big(\sigma_i(t_i)\big)=\phi(X_i)\phi(Y_i)}$, because then
	\[\phi(X_i)\phi(t_i)=\phi(X_i)\phi(Y_i)\phi(X_i)=\phi\big(\sigma_i(t_i)\big)\phi(X_i)\]
	and for $j\neq i$
	\[ \phi(X_i)\phi(t_j)=\phi(X_i)\phi(Y_j)\phi(X_j)=\hat{x}_i\hat{\partial}_j\hat{x}_j=\hat{x}_j\hat{\partial}_j\hat{x}_i=\phi\big(\sigma_i(t_j)\big)\phi(X_i);\]
	likewise for $Y_i$.
	
	Actually, that $\phi(t_i)=\phi(Y_i)\phi(X_i)$ is immediate by definition of $\phi(t_i)$. So, what remains is to prove that
	\begin{equation}\label{eq:pf-phi-1}
		\phi\big(\sigma_i(t_i))=\phi(X_i)\phi(Y_i)
	\end{equation}
	for $i=1,2$.
	
	Now the left-hand side of \eqref{eq:pf-phi-1} equals
	\begin{align}
		\phi(\sigma_i(t_i))&=\phi\bigl(c_i+\hat f_{i1}t_1+\hat f_{i2}t_2\bigr)=c_i+\hat f_{i1}\hat{\partial}_1\hat{x}_1+\hat f_{i2}\hat{\partial}_2\hat{x}_2 \nonumber\\
		&=c_i+\hat f_{i1}\cdot(H_\alpha+2)(H_{\beta+\alpha}+2)\bar{\partial}_1\bar{x}_1+\hat f_{i2}\cdot (H_\alpha+1)(H_{\beta+\alpha}+2)\bar{\partial}_2\bar{x}_2\nonumber\\
		&=c_i+(H_\alpha+i)(H_{\beta+\alpha}+1)\bigl(f_{i1}\bar{\partial}_1\bar{x}_1+f_{i2}\bar{\partial}_2\bar{x}_2\bigr);\label{eq:pf-phi-2}
	\end{align}
	indeed, the right-hand side of \eqref{eq:pf-phi-1} equals
	\[\hat{x}_i\hat{\partial}_i=(H_\alpha+i)(H_{\beta+\alpha}+1)\bar{x}_i\bar{\partial}_i,\]
	which equals \eqref{eq:pf-phi-2}, by \eqref{eq:rel-x1p1}.
	
	At this point we can establish that $\sigma_1$ and $\sigma_2$ actually commute.
	We have for any $b\in B$,
	\[0=(\hat{x}_1\hat{x}_2 -\hat{x}_2\hat{x}_1)b = (\sigma_1(\sigma_2(b))-\sigma_2(\sigma_1(b)))\hat{x}_1\hat{x}_2.\]
	By Corollary \ref{cor:Dsp4-structure:domain}, it follows that $\sigma_1$ and $\sigma_2$ commute.
	
	We have shown all GWA relations \eqref{eq:GWA-rels} are preserved by $\phi$. Therefore, there exists a well-defined $\mathbb{C}$-algebra homomorphism $\phi\colon B(\sigma,t)\to D(\mathfrak{sp}(4))$ as in the statement of Theorem \ref{thm:GWA-realization}. Moreover, $\phi$ is a homomorphism of left $R$-modules.
	By \cite{bavula1992generalized}, when $B=R[t_1,t_2,\dots,t_n]$ for some subring $R$, the monomials $Y_1^aY_2^bX_1^cX_2^d$ ($a,b,c,d\in\mathbb{Z}_{\geq 0}$) form a basis for $B(\sigma,t)$ as a left $R$-module. These are mapped under $\phi$ to the monomials $\hat{\partial}_1^a\hat{\partial}_2^b\hat{x}_1^c\hat{x}_2^d$ which form a left $R$-basis for~$D(\mathfrak{sp}(4))$ by Lemma \ref{lemma:finite-generators}.
	Therefore, $\phi$ is bijective.
\end{proof}

\begin{thebibliography}{99}
% Card-local bibliography fragment for S7-sig25-001-gwa-realization:
% the single cited entry copied verbatim from the pinned source
% thebibliography (source lines 712-714); entry text untouched.
\bibitem{bavula1992generalized}
Bavula V.V., Generalized {W}eyl algebras and their representations,
 \textit{St.~Petersburg Math.~J.} \textbf{4} (1993), 71--92.
\end{thebibliography}
\end{document}
